\documentclass[11pt]{article}

\usepackage[margin=1in]{geometry}
\usepackage{amsmath}
\usepackage{amssymb}
\usepackage{mathtools}
\usepackage{booktabs}
\usepackage{hyperref}
\usepackage{listings}
\usepackage{xcolor}

\hypersetup{
  colorlinks=true,
  linkcolor=blue!50!black,
  urlcolor=blue!50!black,
}

\lstset{
  basicstyle=\ttfamily\small,
  breaklines=true,
  columns=fullflexible,
  keepspaces=true,
  showstringspaces=false,
}

\newcommand{\R}{R}
\newcommand{\fs}{f_{\mathrm{s}}}
\DeclareMathOperator{\sinc}{sinc}

\title{Randomized Multiresolution Band-Energy Search\\
       \large A Sublinear-Time Spectral Peak Finder}
\author{Notes on \texttt{multires\_bandenergy.py}}
\date{\today}

\begin{document}
\maketitle

\begin{abstract}
\texttt{multires\_bandenergy.py} locates a dominant (or local) spectral peak in a
discrete signal while reading a number of time-domain samples that does
\emph{not} grow with the signal length $N$. It never computes an FFT and never
makes a full pass over the data. The trick is to read a small
\emph{autocorrelation sketch} from a handful of random contiguous blocks, then
answer every band-energy question analytically from that sketch via a closed-form
band-integral kernel. A binary (multiresolution) search over frequency bands
then converges to the peak in $O(\log)$ decisions, none of which touch the
signal again. This report derives the underlying mathematics, walks through each
function, and explains why the cost decouples from $N$.
\end{abstract}

\tableofcontents

\section{Problem and Idea}

We are given a real signal $x[0\,..\,N-1]$ sampled at rate $\fs$, and we want the
frequency of a dominant tone (or a local spectral peak near a seed), to some
target resolution $\Delta f$. The textbook route is an FFT: $O(N\log N)$ time and
a full pass over the data. Here we instead want \emph{sublinear} access --- read
$\ll N$ samples --- and we are willing to accept an approximate, Hz-level answer.

The central idea is a divide-and-conquer search in the frequency domain. The
band $[0,\fs/2]$ is split in half; we ask which half carries more energy; we
recurse into the winner. After $d$ levels the band has width $(\fs/2)\,2^{-d}$, so
to reach resolution $\Delta f$ we need only
\begin{equation}
  d \;\approx\; \log_2\!\frac{\fs/2}{\Delta f}
  \label{eq:depth}
\end{equation}
yes/no decisions. The whole game is therefore: \emph{how do we answer
``which half-band has more energy?'' cheaply and reliably}, without an FFT and
without rereading the signal at every level?

The answer is to precompute a small \textbf{autocorrelation sketch} once, and to
turn each band-energy query into a closed-form sum over that sketch.

\section{Mathematical Core}

\subsection{Power spectral density and autocorrelation}

For a wide-sense-stationary process the autocorrelation
\begin{equation}
  \R(\tau) \;=\; \mathbb{E}\,x[n]\,x[n+\tau]
  \;\approx\; \frac{1}{N}\sum_n x[n]\,x[n+\tau]
  \label{eq:acf}
\end{equation}
and the power spectral density (PSD) $S(f)$ are a Fourier-transform pair
(Wiener--Khinchin). With $\tau$ measured in samples and $f$ in Hz,
\begin{equation}
  S(f) \;=\; \sum_{\tau=-\infty}^{\infty} \R(\tau)\,
             e^{-i\,2\pi f \tau/\fs}
       \;=\; \R(0) + 2\sum_{\tau\ge 1}\R(\tau)\cos\!\Big(\frac{2\pi f\tau}{\fs}\Big),
  \label{eq:psd}
\end{equation}
where the second form uses the fact that $\R$ is real and even
($\R(-\tau)=\R(\tau)$), so the imaginary parts cancel and each lag contributes a
cosine.

\subsection{Band energy as a closed form}

The energy in a band $[f_{\mathrm{lo}}, f_{\mathrm{hi}}]$ is the integral of the
PSD over that band:
\begin{equation}
  E[f_{\mathrm{lo}},f_{\mathrm{hi}}]
   \;=\; \int_{f_{\mathrm{lo}}}^{f_{\mathrm{hi}}} S(f)\,df .
\end{equation}
Substituting~\eqref{eq:psd} and exchanging sum and integral,
\begin{equation}
  \boxed{\;
  E[f_{\mathrm{lo}},f_{\mathrm{hi}}]
   \;=\; \R(0)\,(f_{\mathrm{hi}}-f_{\mathrm{lo}})
        \;+\; 2\sum_{\tau\ge 1}\R(\tau)\,g(\tau)
  \;}
  \label{eq:bandenergy}
\end{equation}
where the per-lag \emph{band kernel} is
\begin{equation}
  g(\tau) \;=\; \int_{f_{\mathrm{lo}}}^{f_{\mathrm{hi}}}
                \cos\!\Big(\frac{2\pi f\tau}{\fs}\Big)\,df .
  \label{eq:kerneldef}
\end{equation}
This integral is elementary. For $\tau\ne 0$,
\begin{equation}
  g(\tau) \;=\; \frac{\fs}{2\pi\tau}
    \left[\sin\!\Big(\frac{2\pi f_{\mathrm{hi}}\tau}{\fs}\Big)
        - \sin\!\Big(\frac{2\pi f_{\mathrm{lo}}\tau}{\fs}\Big)\right],
  \qquad
  g(0) \;=\; f_{\mathrm{hi}} - f_{\mathrm{lo}}.
  \label{eq:kernel}
\end{equation}
Equations~\eqref{eq:bandenergy}--\eqref{eq:kernel} are the heart of the module:
\textbf{once $\R(\tau)$ is known for $\tau=0\,..\,M$, the energy in any band is a
finite, closed-form sum --- no further data access, no FFT.} In practice $\R$ is
truncated at a maximum lag $M$, so the sum runs $\tau=1\,..\,M$; this truncation
is exactly what limits the achievable frequency resolution to $\sim \fs/M$ (the
kernel oscillates in $f$ with period $\fs/\tau$, and the fastest term we retain
has period $\fs/M$).

\subsection{Why a contiguous block --- coherent gain}

The reliability of each binary decision comes from a signal-processing fact, not
from sheer sample count. For a tone $A\sin(2\pi f_0 t)$ in white noise of variance
$\sigma^2$, the block autocorrelation is
\begin{equation}
  \R(\tau) \;\approx\; \frac{A^2}{2}\cos\!\Big(\frac{2\pi f_0\tau}{\fs}\Big)
                       \;+\; \sigma^2\,\delta_{\tau,0}.
  \label{eq:tonepluse}
\end{equation}
The tone leaves a \emph{clean cosine at every lag} $\tau=1\,..\,M$, while the
noise contributes \emph{only} at $\tau=0$. When the band in~\eqref{eq:bandenergy}
contains $f_0$, the cosine in~\eqref{eq:tonepluse} and the kernel
$g(\tau)$ line up constructively across all $M$ lags --- an $O(M)$ coherent
build-up. When the band excludes $f_0$, cosine and kernel beat against each other
and the terms cancel. This $O(M)$ contrast between in-band and out-of-band is what
makes each yes/no decision correct.

Crucially this build-up requires \emph{phase-coherent} samples, which a
contiguous block of length $L$ supplies. Scattered point samples (random
frequencies, or random time-pairs $(n,n')$) carry random relative phase, give no
coherent build-up, and leave only variance to average down --- which is
expensive. This is why the module reads \emph{contiguous blocks} rather than
scattered points.

\section{Code Walkthrough}

The module exposes a small set of mostly \emph{pure} functions (in keeping with
the project's preference), built around the sketch.

\subsection{\texttt{ArraySampler} --- the access layer}

A thin wrapper around a NumPy array offering point reads
\verb|sampler(indices)| and contiguous block reads
\verb|sampler.block(start, length)|. The abstraction matters: for genuinely huge
signals the same search code would work against an \texttt{mmap}, disk, network,
or compressed store, because the search only ever asks for a few blocks.

\subsection{\texttt{band\_energy\_kernel} --- the kernel $g(\tau)$ (pure)}

A direct, vectorized implementation of~\eqref{eq:kernel}. It treats $\tau=0$ as a
special case ($g(0)=f_{\mathrm{hi}}-f_{\mathrm{lo}}$) to avoid the $1/\tau$
singularity and applies the closed form to all nonzero lags at once:
\begin{lstlisting}[language=Python]
out[nz] = (fs / (2.0 * np.pi * tnz)) * (
    np.sin(2.0 * np.pi * f_hi * tnz / fs)
    - np.sin(2.0 * np.pi * f_lo * tnz / fs)
)
out[~nz] = f_hi - f_lo
\end{lstlisting}

\subsection{\texttt{autocorrelation\_sketch} --- estimating $\R(\tau)$}

This is the \emph{only} function that reads the signal. It draws \texttt{n\_blocks}
contiguous blocks of length \texttt{block\_len} ($\approx 4M$ by default, so every
lag has ample support), each starting at a uniformly random clamped position, and
for each block accumulates the biased (Bartlett) lag products
\begin{equation}
  \hat\R_{\text{block}}(\tau)
    \;=\; \frac{1}{L}\sum_{n=0}^{L-1-\tau} \tilde x[n]\,\tilde x[n+\tau],
  \qquad \tau = 0\,..\,M,
  \label{eq:biasedacf}
\end{equation}
where $\tilde x$ is the (optionally) mean-removed block. Note the divisor is the
fixed block length $L$ (not the per-lag count $L-\tau$): this is the
\emph{biased} estimator, which tapers the high lags and keeps the kernel sum
in~\eqref{eq:bandenergy} well-behaved (the unbiased version inflates high-lag
variance). The per-block estimates are then averaged over the
\texttt{n\_blocks} blocks to reduce variance. The total signal access is
\begin{equation}
  \texttt{samples\_read} \;=\; \texttt{n\_blocks}\times\texttt{block\_len},
  \label{eq:samples}
\end{equation}
which contains no $N$.

\subsection{\texttt{band\_energy\_from\_sketch} --- evaluating $E$ (pure)}

A literal transcription of~\eqref{eq:bandenergy}: build the lag vector
$\tau=0\,..\,M$, evaluate $g(\tau)$ once, and return
$\R(0)g(0) + 2\sum_{\tau\ge1}\R(\tau)g(\tau)$. It validates
$0\le f_{\mathrm{lo}} < f_{\mathrm{hi}} \le \fs/2$. Because only \emph{relative}
band comparisons drive the search, no absolute normalization is applied.

\subsection{\texttt{multiresolution\_band\_search} --- global search (pure)}

The divide-and-conquer search of Section~2. At each level it splits the current
band $[lo,hi]$ at the midpoint, evaluates the two half-band energies from the
sketch, and keeps the higher-energy half:
\begin{lstlisting}[language=Python]
mid = 0.5 * (lo + hi)
e_left  = band_energy_from_sketch(R, fs, lo,  mid)
e_right = band_energy_from_sketch(R, fs, mid, hi)
if e_left >= e_right:  hi = mid     # keep left
else:                  lo = mid     # keep right
\end{lstlisting}
It stops when the band narrows below \texttt{min\_width\_hz} (default $\fs/M$, the
sketch's own resolution --- searching finer is meaningless) or after
\texttt{max\_depth} levels. The number of levels matches~\eqref{eq:depth}, and
\emph{every level reads zero samples}. It returns the final band, its center, and
optionally a full decision trace. It is reliable when there is a single dominant
band.

\subsection{\texttt{local\_peak\_search} --- hill-climb (pure)}

When a rough location is already known, a local search is cheaper than a global
one (it explores only $\log(\text{radius}/\Delta f)$ levels). Starting from a seed
$f_0$ with an initial step (default $16\times$ resolution), it probes band energy
at the current center and at $\pm\,\text{step}$, moves toward the larger neighbor,
and \emph{halves the step} once the center is the local maximum:
\begin{lstlisting}[language=Python]
if   e_left  > e_center and e_left  >= e_right:  center -= step      # go left
elif e_right > e_center and e_right >  e_left:   center += step      # go right
else:                                            step   *= 0.5       # refine
\end{lstlisting}
Probe bands are clamped to $[0,\fs/2]$. It stops when the step shrinks to the
resolution. This is a 1-D coordinate ascent with geometric step refinement ---
essentially a ternary/golden-style climb on the (sketched) PSD.

\subsection{Convenience wrappers}

\texttt{find\_dominant\_frequency\_sublinear} and
\texttt{find\_local\_peak\_sublinear} tie the pieces together: choose
$M \approx \lceil \fs/\Delta f\rceil$ via \texttt{\_resolution\_to\_max\_lag},
build the sketch once, run the global or local search, and report the recovered
frequency together with bookkeeping (\texttt{samples\_read},
\texttt{sample\_fraction}). They are the user-facing entry points.

\section{Why the Cost Decouples from \texorpdfstring{$N$}{N}}

From~\eqref{eq:samples}, the budget is $\texttt{n\_blocks}\cdot\texttt{block\_len}
\approx \texttt{n\_blocks}\cdot 4M$ --- no $N$. Each factor maps to a real
requirement:

\begin{itemize}
  \item \textbf{Resolution $\to M$.} Resolving to $\Delta f$ needs lags out to
        $M\approx \fs/\Delta f$. A depth-$M$ autocorrelation cannot distinguish
        features finer than $\fs/M$. So samples scale \emph{linearly in target
        resolution}, not in $N$. ($N$ re-enters only if you demand bin-exact
        $\Delta f = \fs/N$, forcing $M\approx N$.)
  \item \textbf{Confidence $\to$ \texttt{n\_blocks}.} More blocks average down the
        variance of $\R(\tau)$. For failure probability $\delta$ over the
        $\sim\log_2(\fs/2/\Delta f)$ decisions (union bound), the block count
        grows like $\log(1/\delta)$. Still no $N$.
  \item \textbf{Contrast $\to$ multiplies \texttt{n\_blocks}.} With $\mathrm{gap}$
        the normalized energy separation between winning and losing half-band,
        reliable decisions need
        \begin{equation}
          \texttt{n\_blocks}\cdot M
            \;=\; O\!\Big(\tfrac{\log(\text{depth}/\delta)}{\mathrm{gap}^2}\Big).
        \end{equation}
        At \emph{fixed SNR} the gap is $N$-independent: as $N$ grows, in-band tone
        energy ($\propto N$) and in-band noise energy ($\propto N$) grow together,
        so their ratio holds. This is the deep reason the budget can stay flat.
  \item \textbf{Search span $\to$ depth.} Logarithmic per~\eqref{eq:depth}, and
        free of samples because every level reuses the one cached sketch.
\end{itemize}

In one line,
\begin{equation}
  \texttt{samples\_read}
    \;\approx\; \texttt{n\_blocks}\cdot\frac{4\fs}{\Delta f}
    \;\approx\; O\!\Big(\frac{\log(1/\delta)}{\mathrm{gap}^2}\cdot\frac{\fs}{\Delta f}\Big).
\end{equation}
The empirical sweep bears this out: a 7130\,Hz tone in noise at
$\Delta f=25$\,Hz, $\texttt{n\_blocks}=4$, reads the \emph{same} 30{,}720 samples
and recovers $\sim$7137\,Hz across $N=10^5,10^6,10^7$ --- a $100\times$ range of
$N$ at identical cost and accuracy.

\section{The Hard Demo: Spectral Leakage}

The \texttt{\_\_main\_\_} block constructs a deliberately adversarial signal: eight
tones of differing strength, all \emph{gated} to a short rectangular burst in the
middle of a one-million-sample record. The abrupt on/off edges convolve each
tone's line with a $\sinc$ in frequency, spreading energy into \emph{sidelobes}
(spectral leakage). The dominant 1500\,Hz tone leaks far enough to partly bury the
weak 2300\,Hz tone beside it.

The demo shows three regimes:
\begin{enumerate}
  \item \textbf{Global search} still picks out the strongest (1500\,Hz) band
        despite the forest of weaker peaks and the leakage between them.
  \item \textbf{Local hill-climb} with a seed inside each tone's own neighborhood
        recovers every peak --- even the masked 2300\,Hz one.
  \item \textbf{Failure mode:} seeding the weak 2300\,Hz tone from farther out
        (2000\,Hz) lets the dominant tone's sidelobe leakage drag the climb into
        the 1500\,Hz peak instead. This honestly illustrates the method's limit:
        local search follows the local energy gradient, and leakage can reshape
        that gradient.
\end{enumerate}

\section{Summary}

The module reduces spectral peak finding to (i) a one-time autocorrelation sketch
read from a few coherent blocks, and (ii) a closed-form band-energy
evaluation~\eqref{eq:bandenergy}--\eqref{eq:kernel} that drives a logarithmic-depth
search reading no further samples. The work is governed by \emph{how finely and
how confidently} you want the peak located --- resolution $\Delta f$, confidence
$\delta$, spectral contrast --- and not by \emph{how long} the signal is. The
honest caveats are that it assumes a dominant (or seed-localizable) peak, that
truncating $\R$ at lag $M$ caps resolution at $\fs/M$, and that strong spectral
leakage can mislead the local climb.

\end{document}
